Recursive bisection for congressional redistricting offers a critical advantage over n-way graph partitioning: \textbf{transparency}. Each district emerges from a hierarchical sequence of binary splits, creating an auditable decision tree that citizens and courts can verify. This contrasts sharply with n-way optimization, which produces all districts simultaneously through iterative refinement—a process that resembles a ``black box'' 	o non-experts.

However, recursive bisection has been assumed 	o suffer from a fundamental weakness: \textbf{sensitivity 	o tree structure}. The sequence of splits—whether 	o divide a state into [3,4] districts or [4,3], whether 	o split the left subgraph as [2,1] or [1,2]—would seemingly affect final district demographics. Prior work on unweighted graph partitioning suggests tree structure matters significantly \citep{karypis1998}.

\subsection{The Tree Structure Question}

Consider South Carolina (k=7 districts, 35.1\% total minority). A balanced binary tree requires three levels:
\begin{itemize}
\item Level 1: Split state into [3,4] or [4,3]
\item Level 2: Split left subgraph into [1,2] or [2,1]; right subgraph into [2,2]
\item Level 3: Split singleton subgraphs into [1] or recurse
\end{itemize}

Traditional recursive bisection uses predetermined trees—the algorithm designer chooses a structure (e.g., [3,4]) before seeing data. This raises three concerns:

\textbf{(1) Arbitrary choices}: Why [3,4] rather than [4,3]? No data-driven justification exists for predetermined structure.

\textbf{(2) Potential suboptimality}: The first split commits 	o grouping districts that may have incompatible demographic goals. If the minority population concentrates in regions best represented by 4 districts, a [3,4] split might separate these regions prematurely.

\textbf{(3) Performance uncertainty}: Without testing alternative trees, we cannot know whether predetermined structures achieve optimal Voting Rights Act (VRA) compliance.

\subsection{Testing Tree Structure Sensitivity}

To test whether tree structure matters for VRA compliance, this paper evaluates three partitioning approaches:

\textbf{(1) Predetermined trees}: Traditional recursive bisection with fixed structure (e.g., [3,4]) chosen before seeing data.

\textbf{(2) Adaptive bisection}: Data-driven tree selection at each split. At every level creating $k_L$ and $k_R$ districts, evaluate both orderings $(k_L, k_R)$ and $(k_R, k_L)$. Choose the ordering that produces higher minority concentration. Recurse with locally optimal structures.

\textbf{(3) N-way partitioning}: Global optimization producing all districts simultaneously through iterative refinement (METIS direct k-way).

All three methods use identical edge-weighting: minority-minority edges receive weight $\alpha = 5$, regular edges weight 1. Edge-weighting creates strong optimization pressure 	o avoid cutting minority clusters \citep{dellaLibera2026nwayRecursive}.

\subsection{Research Questions}

This paper addresses three questions:

\textbf{Q1}: Does tree structure affect VRA compliance outcomes with edge-weighting? Traditional graph partitioning theory suggests different tree structures should produce different results. However, edge-weighting may create such strong optimization signal that tree structure becomes irrelevant.

\textbf{Q2}: Does n-way optimization outperform recursive bisection with edge-weighting? N-way performs global optimization with iterative vertex reassignment, while recursive bisection makes greedy local decisions. Without edge-weighting, n-way typically outperforms recursive methods. But with strong edge-weighting, both may converge 	o the same solution.

\textbf{Q3}: What is the minimum weight factor that produces method-independent results? If edge-weighting does eliminate method differences, at what threshold $\alpha_{crit}$ does convergence occur?

\subsection{Contributions}

\textbf{(1) Empirical discovery}: With edge-weighting at $\alpha = 5$, all three partitioning methods produce \textit{identical} results—same maximum minority percentages, same number of majority-minority districts, same district-level demographics. Tree structure selection and the choice between recursive vs n-way optimization become completely irrelevant.

\textbf{(2) Theoretical framework}: Explains when and why edge-weighting makes method selection irrelevant. Strong edge-weighting ($\alpha \geq 5$) creates optimization signal so powerful that the optimal partition is uniquely determined by geography and population constraints, not algorithmic approach.

\textbf{(3) Comprehensive testing}: Evaluates 27 predetermined tree structures, adaptive bisection, and n-way partitioning across five states with diverse demographics (Alabama, Georgia, Louisiana, Mississippi, South Carolina) and district counts (k = 4 	o 14). All 43 experiments converge 	o identical solutions.

\textbf{(4) Implementation simplification}: Demonstrates that practitioners can use the simplest method (predetermined balanced trees) with confidence. No need for sophisticated adaptive selection or expensive n-way optimization—edge-weighting makes them equivalent.

\subsection{Key Findings}

With edge-weighting at $\alpha = 5$ and minority threshold $\tau = 0.40$, \textbf{all partitioning methods produce identical results}:

\textbf{Alabama (k=7)}: All 6 predetermined trees, adaptive bisection, and n-way partitioning produce 2 majority-minority districts with maximum 50.8\% minority representation. Not merely similar—\textit{identical} down 	o individual district demographics.

\textbf{Georgia (k=14)}: All 13 predetermined trees, adaptive bisection, and n-way produce 6 majority-minority districts with maximum 83.8\%. Every district has identical minority percentage across all methods.

\textbf{Louisiana, Mississippi, South Carolina}: Same pattern—zero variance across methods (p = 1.0, Cohen's d undefined due 	o zero standard deviation).

This finding contradicts our initial hypothesis that adaptive selection would improve over predetermined trees. Instead, we discovered something more fundamental: \textbf{edge-weighting makes method selection irrelevant}. The optimization signal is so strong that greedy recursive decisions, adaptive tree selection, and global n-way optimization all converge 	o the unique partition minimizing weighted edge cuts subject 	o population balance and contiguity constraints.

\textbf{Practical implication}: Use the simplest method (predetermined balanced trees) with confidence. Sophisticated adaptive selection and expensive n-way optimization provide zero benefit when edge-weighting is sufficiently strong.

\subsection{Research Program Context}

This paper addresses a methodological uncertainty that emerged from the broader algorithmic redistricting research program using recursive bisection~\cite{della-libera-recursive-bisection,della-libera-edge-weighted}. While those foundational studies demonstrated that edge-weighted optimization achieves VRA compliance and compactness goals, a critical question remained unanswered: does the choice of tree structure affect outcomes, requiring sophisticated adaptive selection, or can practitioners confidently use simple predetermined trees? This work resolves that question by systematically testing 27 tree structures, adaptive bisection, and n-way optimization across five states, discovering that edge-weighting ($\alpha \geq 5$) makes method selection completely irrelevant—all approaches converge to identical solutions. This finding simplifies implementation for all subsequent work in this program (threshold analysis~\cite{della-libera-threshold-analysis}, temporal stability~\cite{della-libera-temporal-stability}, compactness tradeoffs~\cite{della-libera-compactness-tradeoff}) by establishing that simple balanced recursive bisection is sufficient. Related papers investigate the constraint conflict theory explaining why edge-weighting outperforms multi-constraint optimization~\cite{della-libera-multi-vs-edge}.

\subsection{Paper Organization}

Section 2 provides background on binary tree structures and VRA compliance through edge-weighting. Section 3 describes the adaptive bisection algorithm with pseudocode and complexity analysis (included because we tested it, even though results show it unnecessary). Section 4 presents theoretical framework explaining when edge-weighting produces method-independent results—the signal strength hypothesis and convergence conditions. Section 5 details experimental design testing 27 predetermined trees, adaptive bisection, and n-way partitioning across five states. Section 6 reports results demonstrating complete method equivalence with statistical analysis. Section 7 discusses practical implications—why edge-weighting dominates method choice and how this simplifies implementation. Section 8 concludes with recommendations and future work on weight factor thresholds.
